% GENERATED from content/data/enrichment.json by mxbuild.enrichment; do not edit.
\mxunit{enrichment-lattice-paths-solutions}{Chemins, combinaisons et Pascal}
\mxsubtitle{Première--Terminale · dénombrement \textperiodcentered{} Solutions détaillées \textperiodcentered{} 05}

\begin{exercise}
Calculer $\binom{8}{2}$ et expliquer pourquoi il vaut $\binom{8}{6}$.
\end{exercise}
\begin{solution}
$\binom{8}{2} = 8 \times  \frac{7}{2} = 28$. Choisir deux exclus équivaut à choisir les six retenus.
\end{solution}

\begin{exercise}
Compter les chemins monotones de $\pt{0;0}$ à $\pt{4;3}$.
\end{exercise}
\begin{solution}
Il faut quatre pas à droite parmi sept : $\binom{7}{4} = 35$.
\end{solution}

\begin{exercise}
Parmi ces chemins, combien passent par $\pt{2;1}$ ?
\end{exercise}
\begin{solution}
$\binom{3}{2} \times  \binom{4}{2} = 3 \times  6 = 18$.
\end{solution}

\begin{exercise}
Une équipe de trois personnes est choisie parmi huit. Deux personnes fixées ne peuvent pas être ensemble. Combien d’équipes restent possibles ?
\end{exercise}
\begin{solution}
Il y a $\binom{8}{3} = 56$ équipes. Celles contenant les deux personnes interdites sont au nombre de 6. Il reste 50 équipes.
\end{solution}

\begin{problem}[Chercher, justifier, conclure]
Compter les chemins monotones de $\pt{0;0}$ à $\pt{5;4}$ qui évitent $A\pt{2;1}$ et $B\pt{3;3}$. Commencer par tous les chemins, puis compter ceux passant par A, par B et par les deux. Expliquer pourquoi la dernière catégorie doit être rajoutée.
\end{problem}
\begin{solution}
Total : $\binom{9}{5} = 126$. Par A : $\binom{3}{2}\binom{6}{3} = 60$. Par B : $\binom{6}{3}\binom{3}{2} = 60$. Par les deux : $3 \times  \binom{3}{1} \times  3 = 27$. Les chemins passant par les deux ont été soustraits deux fois ; les rajouter une fois donne $126 - 60 - 60 + 27 = 33$.
\end{solution}

